% fig6-format-decomposition.tex -- Decomposition of TSCG benefit: Format Effect vs Compression Effect
% Included via % fig6-format-decomposition.tex -- Decomposition of TSCG benefit: Format Effect vs Compression Effect
% Included via % fig6-format-decomposition.tex -- Decomposition of TSCG benefit: Format Effect vs Compression Effect
% Included via % fig6-format-decomposition.tex -- Decomposition of TSCG benefit: Format Effect vs Compression Effect
% Included via \input{figures/fig6-format-decomposition} from main.tex
\begin{figure*}[t]
\centering
\begin{tikzpicture}
\begin{axis}[
    ybar,
    bar width=16pt,
    width=\textwidth,
    height=7.5cm,
    ylabel={Accuracy (\%)},
    ylabel style={font=\small},
    xlabel style={font=\small},
    ymin=0,
    ymax=105,
    ytick={0,20,40,60,80,100},
    y tick label style={font=\scriptsize},
    %% Three model groups, each with 3 bars
    xtick={2, 6, 10},
    xticklabels={Claude Sonnet 4, GPT-4o, GPT-5.2},
    x tick label style={font=\small},
    xmin=0, xmax=12,
    nodes near coords,
    nodes near coords style={font=\scriptsize\bfseries, anchor=south, yshift=1pt},
    point meta=explicit symbolic,
    grid=major,
    grid style={dashed, gray!25},
    legend style={
        at={(0.5,0.97)},
        anchor=north,
        legend columns=3,
        font=\small,
        draw=tscggray!40,
        fill=white,
        fill opacity=0.92,
        text opacity=1,
        rounded corners=1pt,
        /tikz/every even column/.append style={column sep=6pt},
    },
    legend cell align={left},
    clip=false,
]

%% ── Natural (native function calling) ──────────────────────────────────────
\addplot[fill=tscggray!55, draw=tscggray!80] coordinates {
    (1, 74.0)  [74.0]
    (5, 55.5)  [55.5]
    (9, 51.9)  [51.9]
};
\addlegendentry{Natural (FC)}

%% ── Natural Text (JSON as text, no compression) ────────────────────────────
\addplot[fill=tscgorange!60, draw=tscgorange!80] coordinates {
    (2, 55.5)  [55.5]
    (6, 45.0)  [45.0]
    (10, 78.3) [78.3]
};
\addlegendentry{Natural Text}

%% ── TSCG (compressed text) ─────────────────────────────────────────────────
\addplot[fill=tscgblue!65, draw=tscgblue!85] coordinates {
    (3, 85.2)  [85.2]
    (7, 56.5)  [56.5]
    (11, 81.6) [81.6]
};
\addlegendentry{TSCG}

%% ── Annotation: GPT-5.2 decomposition (the most dramatic case) ─────────────
%% Format Effect brace: Natural (51.9) -> Natural Text (78.3) = +26.4pp
\draw[decorate, decoration={brace, amplitude=5pt, mirror}, tscgorange, very thick]
    (axis cs:9.25, 51.9) -- (axis cs:9.25, 78.3)
    node[midway, right=7pt, font=\scriptsize, text=tscgorange, align=left]
    {Format\\[-1pt]\textbf{+26.4\,pp}};

%% Compression Effect brace: Natural Text (78.3) -> TSCG (81.6) = +3.3pp
\draw[decorate, decoration={brace, amplitude=4pt}, tscgblue, very thick]
    (axis cs:11.75, 78.3) -- (axis cs:11.75, 81.6)
    node[midway, right=7pt, font=\scriptsize, text=tscgblue, align=left]
    {Compression\\[-1pt]\textbf{+3.3\,pp}};

%% Total effect arrow: Natural (51.9) -> TSCG (81.6) = +29.7pp
\draw[-{Stealth[length=5pt]}, tscggreen, very thick]
    (axis cs:9, 52) -- (axis cs:11, 82)
    node[midway, above=4pt, sloped, font=\scriptsize\bfseries, text=tscggreen]
    {Total: +29.7\,pp};

\end{axis}
\end{tikzpicture}
\caption{Decomposition of TSCG's accuracy benefit into \textbf{Format Effect} (switching from native function calling to uncompressed text) and \textbf{Compression Effect} (applying TSCG compression to text). Scenario~A (16 tools, BFCL). For Claude Sonnet~4, function calling is the superior format (Format~$=$~$-$18.5\,pp), and the entire TSCG benefit comes from compression ({+}29.7\,pp). For GPT-5.2, the pattern reverses dramatically: simply presenting JSON schemas as text yields {+}26.4\,pp, while compression adds only {+}3.3\,pp---indicating that GPT-5.2's native function-calling pathway is the primary bottleneck. GPT-4o shows a smaller version of Claude's pattern.}
\label{fig:format-decomposition}
\end{figure*}
 from main.tex
\begin{figure*}[t]
\centering
\begin{tikzpicture}
\begin{axis}[
    ybar,
    bar width=16pt,
    width=\textwidth,
    height=7.5cm,
    ylabel={Accuracy (\%)},
    ylabel style={font=\small},
    xlabel style={font=\small},
    ymin=0,
    ymax=105,
    ytick={0,20,40,60,80,100},
    y tick label style={font=\scriptsize},
    %% Three model groups, each with 3 bars
    xtick={2, 6, 10},
    xticklabels={Claude Sonnet 4, GPT-4o, GPT-5.2},
    x tick label style={font=\small},
    xmin=0, xmax=12,
    nodes near coords,
    nodes near coords style={font=\scriptsize\bfseries, anchor=south, yshift=1pt},
    point meta=explicit symbolic,
    grid=major,
    grid style={dashed, gray!25},
    legend style={
        at={(0.5,0.97)},
        anchor=north,
        legend columns=3,
        font=\small,
        draw=tscggray!40,
        fill=white,
        fill opacity=0.92,
        text opacity=1,
        rounded corners=1pt,
        /tikz/every even column/.append style={column sep=6pt},
    },
    legend cell align={left},
    clip=false,
]

%% ── Natural (native function calling) ──────────────────────────────────────
\addplot[fill=tscggray!55, draw=tscggray!80] coordinates {
    (1, 74.0)  [74.0]
    (5, 55.5)  [55.5]
    (9, 51.9)  [51.9]
};
\addlegendentry{Natural (FC)}

%% ── Natural Text (JSON as text, no compression) ────────────────────────────
\addplot[fill=tscgorange!60, draw=tscgorange!80] coordinates {
    (2, 55.5)  [55.5]
    (6, 45.0)  [45.0]
    (10, 78.3) [78.3]
};
\addlegendentry{Natural Text}

%% ── TSCG (compressed text) ─────────────────────────────────────────────────
\addplot[fill=tscgblue!65, draw=tscgblue!85] coordinates {
    (3, 85.2)  [85.2]
    (7, 56.5)  [56.5]
    (11, 81.6) [81.6]
};
\addlegendentry{TSCG}

%% ── Annotation: GPT-5.2 decomposition (the most dramatic case) ─────────────
%% Format Effect brace: Natural (51.9) -> Natural Text (78.3) = +26.4pp
\draw[decorate, decoration={brace, amplitude=5pt, mirror}, tscgorange, very thick]
    (axis cs:9.25, 51.9) -- (axis cs:9.25, 78.3)
    node[midway, right=7pt, font=\scriptsize, text=tscgorange, align=left]
    {Format\\[-1pt]\textbf{+26.4\,pp}};

%% Compression Effect brace: Natural Text (78.3) -> TSCG (81.6) = +3.3pp
\draw[decorate, decoration={brace, amplitude=4pt}, tscgblue, very thick]
    (axis cs:11.75, 78.3) -- (axis cs:11.75, 81.6)
    node[midway, right=7pt, font=\scriptsize, text=tscgblue, align=left]
    {Compression\\[-1pt]\textbf{+3.3\,pp}};

%% Total effect arrow: Natural (51.9) -> TSCG (81.6) = +29.7pp
\draw[-{Stealth[length=5pt]}, tscggreen, very thick]
    (axis cs:9, 52) -- (axis cs:11, 82)
    node[midway, above=4pt, sloped, font=\scriptsize\bfseries, text=tscggreen]
    {Total: +29.7\,pp};

\end{axis}
\end{tikzpicture}
\caption{Decomposition of TSCG's accuracy benefit into \textbf{Format Effect} (switching from native function calling to uncompressed text) and \textbf{Compression Effect} (applying TSCG compression to text). Scenario~A (16 tools, BFCL). For Claude Sonnet~4, function calling is the superior format (Format~$=$~$-$18.5\,pp), and the entire TSCG benefit comes from compression ({+}29.7\,pp). For GPT-5.2, the pattern reverses dramatically: simply presenting JSON schemas as text yields {+}26.4\,pp, while compression adds only {+}3.3\,pp---indicating that GPT-5.2's native function-calling pathway is the primary bottleneck. GPT-4o shows a smaller version of Claude's pattern.}
\label{fig:format-decomposition}
\end{figure*}
 from main.tex
\begin{figure*}[t]
\centering
\begin{tikzpicture}
\begin{axis}[
    ybar,
    bar width=16pt,
    width=\textwidth,
    height=7.5cm,
    ylabel={Accuracy (\%)},
    ylabel style={font=\small},
    xlabel style={font=\small},
    ymin=0,
    ymax=105,
    ytick={0,20,40,60,80,100},
    y tick label style={font=\scriptsize},
    %% Three model groups, each with 3 bars
    xtick={2, 6, 10},
    xticklabels={Claude Sonnet 4, GPT-4o, GPT-5.2},
    x tick label style={font=\small},
    xmin=0, xmax=12,
    nodes near coords,
    nodes near coords style={font=\scriptsize\bfseries, anchor=south, yshift=1pt},
    point meta=explicit symbolic,
    grid=major,
    grid style={dashed, gray!25},
    legend style={
        at={(0.5,0.97)},
        anchor=north,
        legend columns=3,
        font=\small,
        draw=tscggray!40,
        fill=white,
        fill opacity=0.92,
        text opacity=1,
        rounded corners=1pt,
        /tikz/every even column/.append style={column sep=6pt},
    },
    legend cell align={left},
    clip=false,
]

%% ── Natural (native function calling) ──────────────────────────────────────
\addplot[fill=tscggray!55, draw=tscggray!80] coordinates {
    (1, 74.0)  [74.0]
    (5, 55.5)  [55.5]
    (9, 51.9)  [51.9]
};
\addlegendentry{Natural (FC)}

%% ── Natural Text (JSON as text, no compression) ────────────────────────────
\addplot[fill=tscgorange!60, draw=tscgorange!80] coordinates {
    (2, 55.5)  [55.5]
    (6, 45.0)  [45.0]
    (10, 78.3) [78.3]
};
\addlegendentry{Natural Text}

%% ── TSCG (compressed text) ─────────────────────────────────────────────────
\addplot[fill=tscgblue!65, draw=tscgblue!85] coordinates {
    (3, 85.2)  [85.2]
    (7, 56.5)  [56.5]
    (11, 81.6) [81.6]
};
\addlegendentry{TSCG}

%% ── Annotation: GPT-5.2 decomposition (the most dramatic case) ─────────────
%% Format Effect brace: Natural (51.9) -> Natural Text (78.3) = +26.4pp
\draw[decorate, decoration={brace, amplitude=5pt, mirror}, tscgorange, very thick]
    (axis cs:9.25, 51.9) -- (axis cs:9.25, 78.3)
    node[midway, right=7pt, font=\scriptsize, text=tscgorange, align=left]
    {Format\\[-1pt]\textbf{+26.4\,pp}};

%% Compression Effect brace: Natural Text (78.3) -> TSCG (81.6) = +3.3pp
\draw[decorate, decoration={brace, amplitude=4pt}, tscgblue, very thick]
    (axis cs:11.75, 78.3) -- (axis cs:11.75, 81.6)
    node[midway, right=7pt, font=\scriptsize, text=tscgblue, align=left]
    {Compression\\[-1pt]\textbf{+3.3\,pp}};

%% Total effect arrow: Natural (51.9) -> TSCG (81.6) = +29.7pp
\draw[-{Stealth[length=5pt]}, tscggreen, very thick]
    (axis cs:9, 52) -- (axis cs:11, 82)
    node[midway, above=4pt, sloped, font=\scriptsize\bfseries, text=tscggreen]
    {Total: +29.7\,pp};

\end{axis}
\end{tikzpicture}
\caption{Decomposition of TSCG's accuracy benefit into \textbf{Format Effect} (switching from native function calling to uncompressed text) and \textbf{Compression Effect} (applying TSCG compression to text). Scenario~A (16 tools, BFCL). For Claude Sonnet~4, function calling is the superior format (Format~$=$~$-$18.5\,pp), and the entire TSCG benefit comes from compression ({+}29.7\,pp). For GPT-5.2, the pattern reverses dramatically: simply presenting JSON schemas as text yields {+}26.4\,pp, while compression adds only {+}3.3\,pp---indicating that GPT-5.2's native function-calling pathway is the primary bottleneck. GPT-4o shows a smaller version of Claude's pattern.}
\label{fig:format-decomposition}
\end{figure*}
 from main.tex
\begin{figure*}[t]
\centering
\begin{tikzpicture}
\begin{axis}[
    ybar,
    bar width=16pt,
    width=\textwidth,
    height=7.5cm,
    ylabel={Accuracy (\%)},
    ylabel style={font=\small},
    xlabel style={font=\small},
    ymin=0,
    ymax=105,
    ytick={0,20,40,60,80,100},
    y tick label style={font=\scriptsize},
    %% Three model groups, each with 3 bars
    xtick={2, 6, 10},
    xticklabels={Claude Sonnet 4, GPT-4o, GPT-5.2},
    x tick label style={font=\small},
    xmin=0, xmax=12,
    nodes near coords,
    nodes near coords style={font=\scriptsize\bfseries, anchor=south, yshift=1pt},
    point meta=explicit symbolic,
    grid=major,
    grid style={dashed, gray!25},
    legend style={
        at={(0.5,0.97)},
        anchor=north,
        legend columns=3,
        font=\small,
        draw=tscggray!40,
        fill=white,
        fill opacity=0.92,
        text opacity=1,
        rounded corners=1pt,
        /tikz/every even column/.append style={column sep=6pt},
    },
    legend cell align={left},
    clip=false,
]

%% ── Natural (native function calling) ──────────────────────────────────────
\addplot[fill=tscggray!55, draw=tscggray!80] coordinates {
    (1, 74.0)  [74.0]
    (5, 55.5)  [55.5]
    (9, 51.9)  [51.9]
};
\addlegendentry{Natural (FC)}

%% ── Natural Text (JSON as text, no compression) ────────────────────────────
\addplot[fill=tscgorange!60, draw=tscgorange!80] coordinates {
    (2, 55.5)  [55.5]
    (6, 45.0)  [45.0]
    (10, 78.3) [78.3]
};
\addlegendentry{Natural Text}

%% ── TSCG (compressed text) ─────────────────────────────────────────────────
\addplot[fill=tscgblue!65, draw=tscgblue!85] coordinates {
    (3, 85.2)  [85.2]
    (7, 56.5)  [56.5]
    (11, 81.6) [81.6]
};
\addlegendentry{TSCG}

%% ── Annotation: GPT-5.2 decomposition (the most dramatic case) ─────────────
%% Format Effect brace: Natural (51.9) -> Natural Text (78.3) = +26.4pp
\draw[decorate, decoration={brace, amplitude=5pt, mirror}, tscgorange, very thick]
    (axis cs:9.25, 51.9) -- (axis cs:9.25, 78.3)
    node[midway, right=7pt, font=\scriptsize, text=tscgorange, align=left]
    {Format\\[-1pt]\textbf{+26.4\,pp}};

%% Compression Effect brace: Natural Text (78.3) -> TSCG (81.6) = +3.3pp
\draw[decorate, decoration={brace, amplitude=4pt}, tscgblue, very thick]
    (axis cs:11.75, 78.3) -- (axis cs:11.75, 81.6)
    node[midway, right=7pt, font=\scriptsize, text=tscgblue, align=left]
    {Compression\\[-1pt]\textbf{+3.3\,pp}};

%% Total effect arrow: Natural (51.9) -> TSCG (81.6) = +29.7pp
\draw[-{Stealth[length=5pt]}, tscggreen, very thick]
    (axis cs:9, 52) -- (axis cs:11, 82)
    node[midway, above=4pt, sloped, font=\scriptsize\bfseries, text=tscggreen]
    {Total: +29.7\,pp};

\end{axis}
\end{tikzpicture}
\caption{Decomposition of TSCG's accuracy benefit into \textbf{Format Effect} (switching from native function calling to uncompressed text) and \textbf{Compression Effect} (applying TSCG compression to text). Scenario~A (16 tools, BFCL). For Claude Sonnet~4, function calling is the superior format (Format~$=$~$-$18.5\,pp), and the entire TSCG benefit comes from compression ({+}29.7\,pp). For GPT-5.2, the pattern reverses dramatically: simply presenting JSON schemas as text yields {+}26.4\,pp, while compression adds only {+}3.3\,pp---indicating that GPT-5.2's native function-calling pathway is the primary bottleneck. GPT-4o shows a smaller version of Claude's pattern.}
\label{fig:format-decomposition}
\end{figure*}
